\begin{table*}[t]
\centering
\caption{Discussion-facing TRACE-BiOpt boundary table.}
\label{tab:trace-biopt-discussion-boundary}
\scriptsize
\setlength{\tabcolsep}{4pt}
\begin{tabularx}{\textwidth}{>{\raggedright\arraybackslash}p{0.14\textwidth}>{\raggedright\arraybackslash}p{0.26\textwidth}>{\raggedright\arraybackslash}p{0.24\textwidth}>{\raggedright\arraybackslash}X}
\toprule
Boundary & What still stands & Explicit non-claim & Reviewer takeaway \\
\midrule
Baseline scope & TRACE-BiOpt is current-best rank 1 on all 9/9 tested rows against the row-wise strongest audited comparator from the 21-method registry, and the external audited comparison-class contract still gives 189/189 Holm-corrected significant wins with no surviving tied or better challenger. & Does not imply dominance over untested baselines, other budgets, or other traffic networks. & Read the paper as an external audited comparison-class contract with a multiplicity-robust all-baseline screen, not as a guarantee against every possible future comparator. \\
Calibrated rerun status & 8/9 current-best rows are already promoted from complete replaceable Stage16 calibrated reruns under the same TRACE-BiOpt objective family. & Rows without completed calibrated reruns remain backed by the Stage15 main evidence chain. & Promoted reruns strengthen weak rows inside one objective family, but only completed reruns replace the Stage15-backed route. \\
Theory and solver & The lower-level MAP model, posterior-risk identity, formal CVaR tail-risk epigraph, validation generalization statement, and searched-neighborhood stationarity certificate are all explicit and audited. & Does not prove exact global optimality, universal MAE improvement, or robustness outside the stated assumptions. & The method is principled bilevel network design, but not an exact global MIP solver or a universal MAE theorem. \\
Mechanism diagnostics & Calibration-sensitive, search-budget-sensitive, and certificate-sensitive regimes are identified inside the same TRACE-BiOpt objective family rather than by switching methods. & These diagnostics explain why current-best rows behave as they do; they are not extra dominance rows or universal monotonicity theorems. & Calibration, search-budget, and certificate diagnostics explain why rows move; they do not create extra dominance claims. \\
Stress evidence & The bounded Stage14 PeMS7\_228 frontier is graph-spectral on 8/9 tested stress slices, with deployment posture changing most under block-missing and narrow failure/drift gaps. & These stress slices remain stress-test-only evidence and do not become TRACE-BiOpt robustness dominance claims. & Deployment stress tests are operational screens, not proof that TRACE-BiOpt dominates under every perturbation family. \\
\bottomrule
\end{tabularx}
\end{table*}
